\section{Conclusion}
\label{sec:conclusion}

We have presented a software workflow for time-varying inflow
parameter estimation in urban-flow simulations, built around an
\ESMDA{} smoother run over a sequence of short, contiguous windows.
Two design choices make the rolling-window construction work.
First, the parameter trajectory of each window is flattened into a
vector that exposes every time point of every physical parameter as
a scalar variable in the augmented state, with an option to
\emph{pin} the window-initial time point so that the Kalman update
preserves the seam value inherited from the previous window.
Second, the prior generator is decoupled from the smoother through
a common \paramTSclass{} interface, which both produces
the cold-start prior of the first window and propagates the
posterior of one window into the prior of the next.  This
decoupling lets us swap among four generative models -- a
GP-with-linear-trend prior, a stationary AR(1) prior, an
Ornstein--Uhlenbeck prior, and a critically-damped AR(2) relaxation
-- without touching the assimilation logic.

The default prior, the critically-damped AR(2) relaxation following
Evensen~\cite{evensen2024}, exposes the latent state $(z, w)$
across windows so that the underlying SDE is itself $C^1$-continuous,
while an exponential relaxation toward the external prior ensures
that, far into a window, the prior spread returns to the
climatological $\Sext$.  The implementation uses the closed-form
discrete-time transition for the critically-damped SDE, which
preserves the stationary covariance exactly without substepping.
Three forward models -- \pylbm{}, \pyudales{}, and \pypalm{} --
are supported through a uniform Python interface, allowing the
truth and the assimilation model to be chosen independently for
identical-twin or cross-solver experiments.

\paragraph{Outlook.}
For the academic article we plan to extend this report along four
axes.
\begin{enumerate}
  \item \emph{Choice of prior model.} A systematic comparison of
        the four prior models on identical-twin experiments with a
        common truth, quantifying the trade-off between smoothness
        at window boundaries, ensemble spread, and posterior RMSE.
  \item \emph{Correlation-length sensitivity.} A sweep of $\ell$ over
        a representative range (e.g.\ $\{60, 120, 200, 400\}\,\mathrm{s}$)
        to characterize the dependence of the recovered parameters
        on the assumed time scale, and the impact of mismatched
        correlation lengths between truth and prior.
  \item \emph{Cross-solver assimilation.} Experiments in which the
        truth is generated by one solver (e.g.\ \pylbm{}) and
        assimilation is performed with another (\pyudales{} or
        \pypalm{}), to characterize the robustness of the recovered
        parameters under model error.
  \item \emph{Joint state-and-parameter assimilation.} The
        \stateparamESMDAclass{} variant of the smoother
        already supports a joint augmented state of inflow
        parameters and a low-dimensional state perturbation; we plan
        to exercise it within the rolling-window pipeline to handle
        residual errors not captured by the parameter time-series
        model alone.
\end{enumerate}
A pinning ablation -- disabling
\code{pin\_initial\_time\_point} from window~$1$ onward and
quantifying the resulting discontinuity at window seams -- will
serve as a diagnostic in support of axes~1 and~2.
